% =============================================================================
% 8.3 GESTIÓN DE CONFIGURACIONES Y SECRETOS
% =============================================================================
\section{Gestión de Configuraciones y Secretos}
\label{sec:secretos}

La gestión de secretos es transversal a todo el sistema y crítica para su seguridad. El
principio rector es \textbf{ningún secreto en el código fuente}: toda credencial llega por
variable de entorno y queda fuera del control de versiones. Esta sección documenta cómo se
externaliza la configuración y cómo se impide el rastreo de secretos en el repositorio.

\subsection{Propiedades de Spring y variables del entorno}
\label{subsec:propiedades}

El backend usa el mecanismo de \textbf{perfiles de Spring}
(capítulo~\ref{ch:backend}, sección~\ref{subsec:config_perfiles},
pág.~\pageref{subsec:config_perfiles}): un \pathbreak{application.yml} base con sobrecargas
\pathbreak{application-local.yml} y \pathbreak{application-prod.yml}. Los valores sensibles
se referencian con la sintaxis \comando{\$\{VARIABLE:defecto\}}, de modo que el fichero
versionado contiene el \textbf{nombre} de la variable, nunca su valor. La
tabla~\ref{tab:secretos} clasifica las credenciales por su naturaleza y ubicación.

\vspace{0.2cm}
\noindent
\renewcommand{\arraystretch}{1.3}
\fontsize{8.5}{10.2}\selectfont\sffamily
\begin{tabularx}{\textwidth}{|L{4.2cm}|C{2.0cm}|Y|}
\hline
\rowcolor{headerTabla}
\sffamily\textbf{Credencial} & \sffamily\textbf{Ámbito} & \sffamily\textbf{Ubicación} \\
\hline
\texttt{SUPABASE\_SERVICE\_ROLE\_KEY} & Privada & Solo backend (entorno). \\ \hline
\rowcolor{grisTecnico}
\texttt{GEMINI\_API\_KEY} & Privada & Solo backend (entorno). \\ \hline
\texttt{supabase.jwt-secret} & Privada & Solo backend (entorno). \\ \hline
\rowcolor{grisTecnico}
\texttt{VITE\_SUPABASE\_ANON\_KEY} & Pública (RLS) & Cliente (\textit{bundle}). \\ \hline
\texttt{VITE\_GOOGLE\_CLIENT\_ID} & Pública (origen) & Cliente (\textit{bundle}). \\ \hline
\end{tabularx}
\label{tab:secretos}

\subsection{Exclusión de secretos del control de versiones}
\label{subsec:gitignore}

El repositorio aplica una política estricta de \textbf{exclusión} mediante
\pathbreak{.gitignore}, cuyo objetivo declarado es ``no subir \textit{build}, basura, datos
de usuario, secretos ni \textit{tooling} de IA''. La figura~\ref{fig:secretos-flujo}
muestra la frontera entre lo versionado y lo excluido.

\vspace{0.3cm}
\begin{figure}[h!]
\centering
\begin{tikzpicture}[node distance=0.55cm and 1.3cm,
    ok/.style={rectangle, rounded corners=3pt, draw=c4Datos!70!black, fill=c4Datos!15,
        font=\sffamily\scriptsize, align=center, minimum width=3.0cm, minimum height=0.7cm, inner sep=3pt},
    no/.style={rectangle, rounded corners=3pt, draw=red!70!black, fill=bgAdvertencia,
        font=\sffamily\scriptsize, align=center, minimum width=3.0cm, minimum height=0.7cm, inner sep=3pt},
    rel/.style={-{Stealth[length=2mm]}, semithick}]
    \node[ok] (a) {\texttt{.env.example}\\\scriptsize (plantilla, sin valores)};
    \node[ok, below=of a] (b) {\texttt{application.yml}\\\scriptsize (solo nombres de variables)};
    \node[no, right=2.2cm of a] (c) {\texttt{.env} · \texttt{.env.*}\\\scriptsize (valores reales)};
    \node[no, below=of c] (d) {\texttt{application-prod.yml}\\\scriptsize (perfiles con secretos)};
    \node[font=\sffamily\scriptsize\bfseries, text=c4Datos!50!black] at (0,1.0) {Versionado};
    \node[font=\sffamily\scriptsize\bfseries, text=red!60!black] at (5.0,1.0) {Excluido (.gitignore)};
\end{tikzpicture}
\caption{Frontera entre artefactos versionados y secretos excluidos del repositorio.}
\label{fig:secretos-flujo}
\end{figure}

La tabla~\ref{tab:gitignore} documenta las reglas de exclusión más relevantes del
\pathbreak{.gitignore} real del backend.

\vspace{0.2cm}
\noindent
\renewcommand{\arraystretch}{1.3}
\fontsize{8.5}{10.2}\selectfont\sffamily
\begin{tabularx}{\textwidth}{|L{4.2cm}|Y|}
\hline
\rowcolor{headerTabla}
\sffamily\textbf{Patrón excluido} & \sffamily\textbf{Motivo} \\
\hline
\texttt{.env} · \texttt{.env.*} (salvo \texttt{.env.example}) & Variables con valores reales; se conserva la plantilla. \\ \hline
\rowcolor{grisTecnico}
\texttt{application-prod.yml} · \texttt{application-secret*} & Perfiles de Spring con credenciales. \\ \hline
\texttt{target/} · \texttt{*.jar} & Artefactos de \textit{build} (regenerables). \\ \hline
\rowcolor{grisTecnico}
Llaves, certificados y \textit{keystores} & Material criptográfico sensible. \\ \hline
\texttt{.claude/} · \texttt{.agents/} & \textit{Tooling} de IA, fuera del producto. \\ \hline
\end{tabularx}
\label{tab:gitignore}

\begin{AlertaSeguridad}
La regla \texttt{.env.*} con la excepción \texttt{!.env.example} es un patrón deliberado:
excluye \textbf{todos} los ficheros de entorno con valores reales pero conserva la
\textbf{plantilla} (sin secretos), de modo que un nuevo desarrollador sepa \textit{qué}
variables debe definir sin que ninguna credencial real toque el repositorio. Esta
disciplina, junto a la separación pública/privada de la tabla~\ref{tab:secretos}, cierra el
modelo de gestión de secretos de la plataforma.
\end{AlertaSeguridad}

\bigskip
Con esto se cierra la documentación de verificación, validación y entornos. El
capítulo~\ref{ch:anexos} (pág.~\pageref{ch:anexos}) recopila los anexos y el registro
documental que complementan este manual.
